\chapter{The Bakery Event Corpus}
\label{ch:bakery}

\index{XY plot}
\index{profit!derived from events}
\index{waste!derived from events}
The bakery is the quantitative worked example. Its source is not a table of
percentages copied into this chapter. It is the \BakeryEvents{}-line
\texttt{examples/bakery-events.jsonl} stream. The reducer reads those events and
generates both \texttt{examples/bakery-stats.json} and the coordinates plotted
below. Changing the stream changes the book.

\section{The grammar represented by events}

An order begins with \texttt{OrderPlaced}. Payment and baking fork from that
event, and \texttt{OrderAccepted} or \texttt{OrderRejected} joins their result
events. An accepted order either resolves to customer cancellation or continues
through courier assignment to delivery or delivery failure. The file also
contains \texttt{ShiftPaid} events, one per employee shift. Thus operational and
accounting facts share the same event algebra without being the same task.

For each calendar-day bucket $d$, the reducer computes

\begin{align*}
\operatorname{grossSales}_d
  &= \sum_{e.kind=\mathrm{PaymentAuthorized}} e.amount,\\
\operatorname{refunds}_d
  &= \sum_e e.refund,\\
\operatorname{netRevenue}_d
  &= \operatorname{grossSales}_d-\operatorname{refunds}_d,\\
\operatorname{ingredientCost}_d
  &= \sum_{e.kind\in\{\mathrm{BakeCompleted},\mathrm{BakeFailed}\}} e.ingredientCost,\\
\operatorname{employeePay}_d
  &= \sum_{e.kind=\mathrm{ShiftPaid}} e.pay,\\
\operatorname{profit}_d
  &= \operatorname{netRevenue}_d-\operatorname{ingredientCost}_d
     -\operatorname{employeePay}_d.
\end{align*}

Waste cost is already part of ingredient cost, so subtracting it again would
double count. It is tracked separately to explain lost margin. This modeled
profit excludes rent, utilities, tax, equipment depreciation, and any other cost
for which the stream has no event. Adding such events extends the equation rather
than changing the meaning of existing fields.

\section{Money by day is the primary XY plot}

\begin{figure}[htbp]
\centering
\begin{tikzpicture}
\begin{axis}[
  width=.96\textwidth,height=7.2cm,
  xlabel={day},ylabel={dollars},
  grid=major,legend pos=south west,
  scaled ticks=false]
  \addplot[color=leanblue,mark=none,thick] coordinates {\BakeryDailyRevenueCoordinates};
  \addlegendentry{net revenue}
  \addplot[color=leangreen,mark=*,mark size=1.2pt,thick] coordinates {\BakeryDailyProfitCoordinates};
  \addlegendentry{profit after ingredients and pay}
  \addplot[color=orange!85!black,mark=none,dashed,thick] coordinates {\BakeryDailyPayCoordinates};
  \addlegendentry{employee pay}
\end{axis}
\end{tikzpicture}
\caption{Daily bakery economics reduced from events. Day 35 is a partial day in
the corpus, which explains its lower revenue and profit.}
\label{fig:bakery-daily-money}
\end{figure}

The employee-pay line is not a constant assumption. Each \texttt{ShiftPaid}
event carries employee, role, minutes, hourly rate, and calculated pay. This lets
the same chart answer counterfactual questions after a staffing event stream is
changed.

\section{Waste explains part of the margin}

\begin{figure}[htbp]
\centering
\begin{minipage}{.48\textwidth}
\centering
\begin{tikzpicture}
\begin{axis}[
  width=\linewidth,height=6cm,
  xlabel={day},ylabel={waste cost (dollars)},grid=major,
  scaled ticks=false]
  \addplot[color=red!75!black,mark=none,thick]
    coordinates {\BakeryDailyWasteCostCoordinates};
\end{axis}
\end{tikzpicture}
\end{minipage}\hfill
\begin{minipage}{.48\textwidth}
\centering
\begin{tikzpicture}
\begin{axis}[
  width=\linewidth,height=6cm,
  xlabel={day},ylabel={wasted units},grid=major,
  scaled ticks=false]
  \addplot[color=orange!90!black,mark=square*,mark size=1pt]
    coordinates {\BakeryDailyWasteUnitsCoordinates};
\end{axis}
\end{tikzpicture}
\end{minipage}
\caption{Waste cost and units by day, derived from failed bakes, cancellations,
and failed deliveries in the same event stream.}
\label{fig:bakery-waste}
\end{figure}

Cost and units are deliberately separate. A wasted sourdough loaf and a wasted
croissant do not have the same economic effect. Grouping the same events by
product, outcome, employee shift, or hour of day produces further XY datasets
without changing the scenario.

\section{XY plots as questions over scenarios}

The XY chart is not restricted to time on the horizontal axis. The reducer also
generates daily order count against daily profit:

\begin{figure}[htbp]
\centering
\begin{tikzpicture}
\begin{axis}[
  width=.72\textwidth,height=6.5cm,
  xlabel={orders placed in day},ylabel={profit (dollars)},grid=major,
  scaled ticks=false]
  \addplot[only marks,color=leanblue,mark=*]
    coordinates {\BakeryOrdersProfitCoordinates};
\end{axis}
\end{tikzpicture}
\caption{Daily order volume against profit. Each point is one day reduced from
the event corpus.}
\label{fig:bakery-orders-profit}
\end{figure}

The same mechanism can measure many scenarios: offered load against latency,
staffing hours against profit, price against cancellation probability, queue
capacity against completed orders, or waste against bake batch size. In each
case the axes are named reductions over event streams. A parameter sweep produces
several streams; observed operation supplies streams from different days. The
chart renderer need not know which kind it received.

The corpus currently reports \BakeryDeliveredRatio{} delivered orders, median
terminal latency \BakeryPFiftyLatency{} ms, 95th-percentile latency
\BakeryPNinetyFiveLatency{} ms, and a maximum of \BakeryMaxActive{} active orders.
These values are generated from the checked-in stream and serve as the default
probability and performance inputs for the bakery model.

\section{AI rendering from durable requirements to pseudocode}

The grammar is useful for more than recognizing a saved corpus. It generates
legal traces, so the same durable specification can drive software generation.
The inputs are files, not the conversation that produced them. Let $R$ be the
complete current \texttt{Requirements.lean}, $I$ the complete current
\texttt{Implementation.lean}, $P$ the protobuf schema, and $C_k$ an optional
existing implementation. Code generation has two forms:

\[
  (R,I,P) \longrightarrow C_0,
  \qquad
  (R,I,P,C_k) \longrightarrow C_{k+1}.
\]

The first creates software from scratch. The second revises existing software
against the same stable inputs, retaining conforming code and replacing code
that conflicts. Prior chat messages are deliberately absent from both
functions. If a necessary decision is not present in $R$, $I$, or $P$, the
generator reports the missing durable fact instead of recovering it from a long
conversation.

For the bakery, an AI rendering of those inputs can first produce implementation-
neutral pseudocode. This is a review artifact before a target-language backend
uses filenames, transports, queue APIs, and protobuf bindings from $I$:

\begin{lstlisting}[language={}]
type Event = {
  id, prior[], session, task, srcInstance, dstInstance,
  kind, payload, timeAt, joinRequired?, joinTotal?, joinSelected[]
}

// @requirement bakery.task.fulfill_order
// Justification: realize the legal order-fulfillment production.
function handleOrderPlaced(placed):
  // @requirement bakery.wire.round_trip
  require decode(encode(placed)) == placed
  // @requirement bakery.resource.max_in_flight
  key = (placed.session, placed.task)
  reserveInFlight(StorefrontInstance, key)

  // @requirement bakery.grammar.payment_parallel_bake
  paymentRequest = emit PaymentRequested(
    prior=[placed.id], values={price_cents})
  bakeRequest = emit BakeRequested(
    prior=[placed.id], values={product, quantity})

  // Grammar parallelism: either result may arrive first.
  sendBlocking(Payment.outbound, encode(paymentRequest))
  sendBlocking(Bakery.outbound, encode(bakeRequest))

  // @requirement bakery.join.payment_and_bake
  paymentResult = receiveFor(key, PaymentResult)
  bakeResult = receiveFor(key, BakeResult)

  // Ordinary all-of join: two of two branches are required.
  decision = emit OrderAccepted or OrderRejected(
    prior=[paymentResult.id, bakeResult.id, placed.id],
    joinRequired=2, joinTotal=2,
    joinSelected=[paymentResult.id, bakeResult.id],
    timeAt=clock.now())

  if decision is OrderRejected:
    releaseInFlight(StorefrontInstance, key)
    return

  if probabilisticUserChoice(Cancel, Continue) is Cancel:
    emit OrderCancelled(prior=[decision.id, placed.id])
  else:
    requested = emit DeliveryRequested(prior=[decision.id])
    assigned = await CourierAssigned(backpointing=requested.id)
    emit OrderDelivered or DeliveryFailed(
      prior=[assigned.id, placed.id], timeAt=clock.now())

  releaseInFlight(StorefrontInstance, key)

// @requirement bakery.metrics.event_characterization
// Justification: derive the required observable metrics from events.
function reduce(events):
  assert uniqueIds(events)
  assert acyclicPriorGraph(events)
  assert everyEndBackpointsToItsStart(events)
  assert everyJoinHasValidThresholdEvidence(events)

  latency = end.timeAt - start.timeAt
  clientRate = sum(bytesMoved) / sum(latency)
  serverRate = sum(bytesMoved) /
               (max(end.timeAt) - min(start.timeAt))
  outage = sum(recovered.timeAt - unavailable.timeAt) /
           (replicas * observationWindowMs)
  mttr = sum(recovered.timeAt - unavailable.timeAt) / incidents

  emit interactionDiagramPerScenario(events)
  emit stateMachinePerScenario(events)
  emit xyLines(latency, throughput, profit, waste, outage, mttr)
  emit pieHistograms(outcomes, products, failures)

// @requirement bakery.conformance.legal_trace_bytes
// Justification: generated code must preserve and recognize legal traces.
function conformanceTest(trace):
  bytes = protobufEncode(trace)
  replay = protobufDecode(bytes)
  require replay == trace
  require grammarAccepts(replay)
  require stateMachinesAccept(replay)
  require resourceContractsHold(replay)
\end{lstlisting}

The pseudocode does not become a second source of requirements. It is discarded
and regenerated when $R$ or $I$ changes. Generated source code is likewise a
projection, except that revision mode accepts the current code as an additional
input so hand-written conforming sections need not be destroyed. In both modes,
the grammar supplies positive tests by generating legal traces and negative
tests by rejecting mutations outside its language. Byte round trips then check
that the implementation and its transport bindings preserve the observable
scenario.

Every generated code unit carries a machine-readable requirement reference and
a short justification beside the code it explains. The reference points into
the durable requirement namespace, not to a chat turn. Generation also emits a
traceability manifest mapping requirement IDs to source files, symbols, and
tests. An implementation mapping without a requirement reference is an orphan
and fails implementation-plan validation; generated-source validation must apply
the same rule to emitted symbols. Several code units may cite one requirement, and one code unit may
cite several requirements; the mapping is many-to-many.
